\documentclass[11pt,a4paper]{ltjsarticle}
\usepackage[luatex,pdfencoding=auto]{hyperref}
\usepackage{luatexja-fontspec}
\usepackage{geometry}
\geometry{margin=25mm}

% フォント設定（標準的な設定に修正）
\setmainjfont{HaranoAjiMincho-Regular}
\setsansjfont{HaranoAjiGothic-Medium}

% タイトル周りの調整（手動）
\title{\textsf{\huge 自己紹介および実績報告書}}
\author{\large 高久}
\date{\today}

\begin{document}

\maketitle

\section*{概要}
私は「高度な技術実装」と「社会課題の解決」の両立を信条とするエンジニアです。高専在学中、自社プロジェクト「MEKOU」を立ち上げ、分散型AIインフラの開発や介護ロボティクスへの応用に取り組んできました。

\section*{主な実績とスキルセット}
\begin{itemize}
    \item EJPT
    \item DCON 2026 本選出場
    \item 様々なプロジェクト開発、共同開発ゲーム開発なども経験
    \item 介護職員初任者研修の資格取得中
    \item CTFやハッカソンなどのコンテスト参加経験も豊富
    \item Rust, C系、 JAVA, TSなどのプログラミング言語、AI関連技術、SLAM、ロボティクスなど幅広い技術スタックを保有
    \item 電子工作が趣味の趣味から、EMGなどのプロジェクトも経験
    \item CGも様々な場面で活用することができました。
\end{itemize}

\section{志望分野}
第一志望：情報経営システム、第二志望：電気電子情報工学

\section*{志望動機(第一志望の理由)}
私は、現在の目標として生体信号処理を用いた研究を希望します。
今のように開発に没頭することになったのは個人プロジェクトで行ったEMGの開発でした。当初はEEGの代替として着手しましたが、この経験を通じ、生体信号処理の可能性に魅了されました。
貴校のインターンシップにて、その整った研究環境、学術的な真摯さに触れ、ここで研究したいと思いました。

\section{大学院進学の希望}
    あり

\section{学生生活の抱負}
ある程度苦手なことでも、ホームステイや、documentなどのために英語なども克服しようとしていたり、苦手なことも克服していきたいと思っています。
さらに、貴校の研究室の活動や、プロジェクトを通じて、技術スキルをさらに磨きをかけていきたいと存じます。
また、起業も視野に入れ、常にその心がけをもって、多角的な視点を養う所存です。

\section{将来に対する抱負}
長岡技術科学大学の「実務訓練」および「フロンティアコース」を通じ、現在取り組んでいるMEKOUプロジェクトを学術的・産業的に昇華させたいと考えています。情報、電子、クライアントサイドの知見を融合させ、介護現場などで機能するインフラを構築し、社会へ貢献を目指したいと思います。


\section{自分が優れていると思うこと}
私の強みは、cybersecurityの学習で培った低レイヤーから物事を構築する視点です。
これにより、一般科目においても、理解が深まりました。
高専での学業を疎かにせず、個人開発においてハードウェアからソフトウェア、さらにはビジュアル表現（CG）までを垂直統合的に扱ってきました。物事を根本（低レイヤー）から理解しようとする姿勢が、複雑なシステム設計における課題解決能力に繋がっていると自負しています。


\section{TOEICのスコア}
\end{document}